% TOP_PROBLEM: {"id": 30007101, "title": "Power saving for prime-intersective polynomial differences", "category_id": 2, "category": {"id": 2, "name": "combinatorics", "display_name": "Combinatorics", "description": "Counting problems, graph theory, discrete structures.", "slug": "combinatorics", "order_index": 2, "created_at": "2026-07-31T15:26:25.671Z"}, "status": "open"}
% FIRST_POSTED: 2026-10-06
% ATTEMPT_STATUS: solved
% OPENAI_CLAIM: {"schema_version":1,"source_repo":"https://github.com/openai/math","source_commit":"adc7f1241b42e322a6451854ab7e4b4c146bf78a","release_date":"2026-10-06","match":"full","resolution":"proved","independently_reviewed":false,"second_pass":"agreed","families":[{"family":"182","title":"Power savings for intersective polynomial differences and prime arguments","subject":"Combinatorics","manuscripts":[{"title":"A Power Saving for Polynomial Differences at Prime Arguments","pdf_path":"preprints/A-Power-Saving-for-Polynomial-Differences-at-Prime-Arguments-October-5-2026/prime-argument-polynomial-differences.pdf","readme_path":"preprints/A-Power-Saving-for-Polynomial-Differences-at-Prime-Arguments-October-5-2026/README.md","date":"2026-10-05","theorem_ref":"preprints/A-Power-Saving-for-Polynomial-Differences-at-Prime-Arguments-October-5-2026/build/sections/introduction.tex#thm:main"}],"lean":{"doc_path":"lean/docs/182.md","comparators":[],"declarations":[],"covers_main_theorem":false},"reasoning_trace":null}]}
\documentclass[11pt]{article}
\usepackage[margin=25mm]{geometry}
\usepackage{fontspec}
\setmainfont{Latin Modern Roman}
\usepackage{amsmath,amssymb,mathtools}
\usepackage{unicode-math}
\setmathfont{Latin Modern Math}
\usepackage{enumitem}
\usepackage{xurl,hyperref,bookmark}
\hypersetup{colorlinks=true,urlcolor=blue}
\setcounter{secnumdepth}{0}
\setlength{\parindent}{0pt}
\setlength{\parskip}{5pt}
\setlength{\emergencystretch}{3em}
\begin{document}
\section*{Power saving for prime-intersective polynomial differences}

\subsection{OpenAI's claim}
\textit{Polynomial differences at prime arguments}: For a fixed integer polynomial with positive leading term and a unit root modulo every modulus, the paper claims a fixed power saving for sets avoiding its nonzero values at prime inputs as differences. Both the exponent and constant may depend on the polynomial.

\subsection{Outline of the argument}
\textit{Polynomial differences at prime arguments}: The progression-compatible polynomial-difference method is adapted to a distribution supported on prime arguments. Root normalizations retain the required coprimality through successive arithmetic progressions. Prime-distribution estimates from the companion Dirichlet zero-free result supply the local comparison needed to contract the density-dominating multilinear quantity.

\subsection{Scope relative to this problem}
\textit{Polynomial differences at prime arguments}: Prime-intersective means a root r coprime to q for every q, exactly the unit-root condition. The manuscript’s result quantifies over every such positive-leading h of degree>=2 and every positive N, with c\_h>0,C\_h and avoidance of nonzero h(p) differences at rational prime inputs. Its exponent may depend on h.

With the hypotheses and conventions above, the claimed result addresses the whole mathematical assertion on this page.

\subsection{Formal verification}
Family 182 describes Lean coverage for the square-difference power saving; the general polynomial and prime-input estimates are separate.

\textit{Polynomial differences at prime arguments}: No manuscript-specific Comparator is listed in the pinned inventory.

The documented selection does not cover the entire manuscript assertion relevant here.

This site has not run the Lean development. These links describe OpenAI’s selected formal statements and should be read with their stated scope.

\subsection{Status caveat}
OpenAI reports that an internal model produced the manuscripts. Their arguments have not been independently verified as proofs by this site. Any portion outside the documented formal scope may contain mistakes, and the catalogue label reports the announced mathematical claim.

\subsection{Sources}
\begin{itemize}
\item OpenAI, \emph{A Power Saving for Polynomial Differences at Prime Arguments}, 2026-10-05; citation key \texttt{OAI:A-Power-Saving-for-Polynomial-Differences-at-Prime-Arguments-October-5-2026}. \href{https://github.com/openai/math/blob/adc7f1241b42e322a6451854ab7e4b4c146bf78a/preprints/A-Power-Saving-for-Polynomial-Differences-at-Prime-Arguments-October-5-2026/prime-argument-polynomial-differences.pdf}{PDF (GitHub)}; \href{https://github.com/openai/math/raw/adc7f1241b42e322a6451854ab7e4b4c146bf78a/preprints/A-Power-Saving-for-Polynomial-Differences-at-Prime-Arguments-October-5-2026/prime-argument-polynomial-differences.pdf}{PDF (direct)}; \href{https://github.com/openai/math/blob/adc7f1241b42e322a6451854ab7e4b4c146bf78a/preprints/A-Power-Saving-for-Polynomial-Differences-at-Prime-Arguments-October-5-2026/README.md}{Manuscript page and citation}; \href{https://github.com/openai/math/blob/main/preprints/A-Power-Saving-for-Polynomial-Differences-at-Prime-Arguments-October-5-2026/prime-argument-polynomial-differences.pdf}{Latest version}.
\item Family 182: \href{https://github.com/openai/math/blob/adc7f1241b42e322a6451854ab7e4b4c146bf78a/lean/docs/182.md}{Lean scope}.
\item \href{https://github.com/openai/math/blob/adc7f1241b42e322a6451854ab7e4b4c146bf78a/CONTENTS.md}{OpenAI catalogue entry}.
\end{itemize}
\end{document}
